\usepackage{xfp}
\makeatletter

\usepackage{xcolor}
\iftemplate
  \newcommand{\note}[1]{\textcolor{blue}{\textit{#1}}}
\else
  \newcommand{\note}[1]{}
\fi

% ============================================================
% CONFIGURATION
% ============================================================

% Total number of MTs
\def\MoM@mtcount{5}

% Number of questions per MT
\expandafter\def\csname MT@numq@1\endcsname{7}
\expandafter\def\csname MT@numq@2\endcsname{7}
\expandafter\def\csname MT@numq@3\endcsname{6}
\expandafter\def\csname MT@numq@4\endcsname{4}
\expandafter\def\csname MT@numq@5\endcsname{5}

% ============================================================
% NC / PC / FC  mappings
% ============================================================

% Numeric values (used for the radar average)
\expandafter\def\csname MT@val@NC\endcsname{0}
\expandafter\def\csname MT@val@PC\endcsname{0.5}
\expandafter\def\csname MT@val@FC\endcsname{1}

% Full text labels (used in the per-MT score table)
\expandafter\def\csname MT@text@NC\endcsname{Not covered}
\expandafter\def\csname MT@text@PC\endcsname{Partially covered}
\expandafter\def\csname MT@text@FC\endcsname{Fully covered}

% \MT@MapVal{NC|PC|FC} -> 0 / 0.5 / 1
\newcommand{\MT@MapVal}[1]{%
  \ifcsname MT@val@#1\endcsname
    \csname MT@val@#1\endcsname
  \else
    0%
  \fi
}

% \MT@MapText{NC|PC|FC} -> "Not covered" / "Partially covered" / "Fully covered"
\newcommand{\MT@MapText}[1]{%
  \ifcsname MT@text@#1\endcsname
    \csname MT@text@#1\endcsname
  \else
    \textit{(unset)}%
  \fi
}

% ============================================================
% INTERNALS 
% ============================================================

% \MTNumQ{n} -- expands to the number of questions for MT n
\newcommand{\MTNumQ}[1]{\csname MT@numq@#1\endcsname}

% \MT{<mt>}{<q>}{<coverage>}
% - if <coverage> is non-empty: store the symbolic value (NC/PC/FC)
% - if <coverage> is empty: retrieve it (undefined entries read as NC)
\newcommand{\MT}[3]{%
  \if\relax\detokenize{#3}\relax
    \ifcsname MT@#1@#2\endcsname
      \csname MT@#1@#2\endcsname
    \else
      NC%
    \fi
  \else
    \expandafter\gdef\csname MT@#1@#2\endcsname{#3}%
  \fi
}

% --- Loop counters (two to avoid conflicts in nested loops) ---
\newcount\MT@ctrA
\newcount\MT@ctrB

% ============================================================
% PER-MT AVERAGE (for radar)
% ============================================================

% \MT@BuildAvgExpr{mt} -- builds \MT@avgexpr as an fpeval-ready expression
%   that evaluates to the average mapped value for MT <mt>.
\newcommand{\MT@BuildAvgExpr}[1]{%
  \global\MT@ctrB=0\relax
  \gdef\MT@avgexpr{0}%
  \loop
    \global\advance\MT@ctrB by 1
    \ifnum\MT@ctrB<\numexpr\MTNumQ{#1}+1\relax
      \xdef\MT@avgexpr{%
        \unexpanded\expandafter{\MT@avgexpr}%
        +\noexpand\MT@MapVal{\noexpand\MT{#1}{\the\MT@ctrB}{}}}%
  \repeat
  % divide by number of questions
  \xdef\MT@avgexpr{(\unexpanded\expandafter{\MT@avgexpr})/\MTNumQ{#1}}%
}

% ============================================================
% TABLE ROW BUILDER (3 columns: MT | Question | Coverage)
% ============================================================

\newcommand{\MT@Row}[2]{%
  & \textbf{#2} & \textsf{\MT@MapText{\MT{#1}{#2}{}}} \\
}

\newcommand{\MT@BuildRows}[1]{%
  \global\MT@ctrB=1\relax
  \gdef\MT@builtrows{}%
  \loop
    \global\advance\MT@ctrB by 1
    \ifnum\MT@ctrB<\numexpr\MTNumQ{#1}+1\relax
      \xdef\MT@builtrows{%
        \unexpanded\expandafter{\MT@builtrows}%
        \noexpand\MT@Row{#1}{\the\MT@ctrB}}%
  \repeat
}

% Full multirow block for MT #1
\newcommand{\MT@Block}[1]{%
  \MT@BuildRows{#1}%
  \multirow{\MTNumQ{#1}}{4em}{\textsf{\textbf{MT.#1}}}%
  & \textbf{1} & \textsf{\MT@MapText{\MT{#1}{1}{}}} \\
  \MT@builtrows
  \hline
}

% ============================================================
% RADAR BUILDER
% ============================================================

% Build radar axes + coverage path covering ALL configured MTs
% (1..\MoM@mtcount), regardless of whether \MTScoreTable was called
% for them.  Omitted MTs simply contribute an average of 0 (NC).
%
% \MT@radaraxes -- TikZ \draw / \node commands for the spokes + labels
% \MT@radarpath -- sequence of  (<angle>:<r>) --  coordinates

% Per-MT step.  Lives outside the outer \loop so it can safely use its
% own \loop (via \MT@BuildAvgExpr) -- plain TeX's \loop\repeat is not
% nestable, so we collect a list of \MT@RadarStep calls during the
% outer loop and execute them afterwards.
\newcommand{\MT@RadarStep}[1]{%
  % angle = 90 - 360*(i-1)/N  (clockwise from top, first axis on top)
  \edef\MT@angle{\fpeval{90 - 360*(#1-1)/\MoM@mtcount}}%
  \MT@BuildAvgExpr{#1}%
  \edef\MT@avg{\fpeval{\MT@avgexpr}}%
  \xdef\MT@radaraxes{%
    \unexpanded\expandafter{\MT@radaraxes}%
    \noexpand\draw[gray!60] (0,0) -- (\MT@angle:1);
    \noexpand\node[font=\noexpand\small\noexpand\sffamily]
      at (\MT@angle:1.18) {MT.#1};
  }%
  \xdef\MT@radarpath{%
    \unexpanded\expandafter{\MT@radarpath}%
    (\MT@angle:\MT@avg) --
  }%
  \xdef\MT@radarframe{%
    \unexpanded\expandafter{\MT@radarframe}%
    (\MT@angle:1) --
  }%
}

\newcommand{\MT@BuildRadarPath}{%
  % Phase 1: queue one \MT@RadarStep call per MT (no nested \loop yet).
  \gdef\MT@radarcalls{}%
  \global\MT@ctrA=0\relax
  \loop
    \global\advance\MT@ctrA by 1
    \ifnum\MT@ctrA<\numexpr\MoM@mtcount+1\relax
      \xdef\MT@radarcalls{%
        \unexpanded\expandafter{\MT@radarcalls}%
        \noexpand\MT@RadarStep{\the\MT@ctrA}}%
  \repeat
  % Phase 2: execute the queued calls (each may safely use \loop).
  \gdef\MT@radaraxes{}%
  \gdef\MT@radarpath{}%
  \gdef\MT@radarframe{}%
  \MT@radarcalls
}

% ============================================================
% PUBLIC INTERFACE
% ============================================================

% \MTScoreTable{n} -- insert a coverage table for MT n
\newcommand{\MTScoreTable}[1]{%
  \medskip
  \begin{center}
  \begin{tabular}{ |c|c|c| }
  \cline{2-3}
  \multicolumn{1}{c|}{} & \textbf{Question} & \textbf{Coverage} \\
  \hline
  \MT@Block{#1}%
  \end{tabular}
  \end{center}
}

% \MTRadar -- draw a radar / spider diagram of per-MT average coverage
\newcommand{\MTRadar}{%
  \MT@BuildRadarPath
  \begin{center}
  \begin{tikzpicture}[scale=3]
    % Concentric reference polygons (NC=0, PC=0.5, FC=1 levels)
    \foreach \r in {0.25,0.5,0.75,1.0}{%
      \begin{scope}[scale=\r]
        \draw[gray!25, thin] \MT@radarframe cycle;
      \end{scope}%
    }
    % Spokes and MT labels
    \MT@radaraxes
    % Coverage polygon
    \draw[gray!60!black, thick, fill=gray!40, fill opacity=0.55]
      \MT@radarpath cycle;
  \end{tikzpicture}
  \end{center}
}

\makeatother
